% Problem record op_5ea7cb9abb281829. This TeX file is derived from
% database/problems_json/op_5ea7cb9abb281829.json by scripts/sync-tex.mjs; edit the JSON record
% and rerun the script rather than editing this file.
\section{Optimal dimension of seven-qubit distance-two codes}
\label{sec:op_5ea7cb9abb281829}

\paragraph{Problem.}
For an exact seven-qubit code that detects every single-qubit error, what is
the largest possible code-space dimension: is $K_{\max}^{(2)}(7,2)$ equal to
$24$, $25$, or $26$?  Let $\mathcal P_7:=\{I,X,Y,Z\}^{\otimes7}$ be the Pauli
basis, where the weight $\operatorname{wt}(E)$ counts nonidentity tensor
factors, and define
\begin{equation}
  K_{\max}^{(2)}(7,2):=\max\Bigl\{\operatorname{rank}P:\
  P=P^\dagger=P^2\ \text{on}\ (\mathbb C^2)^{\otimes7},\
  PEP=c_EP\ \text{for every}\ E\in\mathcal P_7\ \text{with}\
  \operatorname{wt}(E)=1\Bigr\},
  \label{eq:5ea7-kmax}
\end{equation}
where each $c_E$ is a scalar depending only on $E$.  The maximum in
Eq.~\eqref{eq:5ea7-kmax} is over all projectors, including degenerate, impure,
and non-CWS codes.  A distance-two code detects an arbitrary single-qubit error
and corrects an erasure at a known location; it need not correct an arbitrary
single-qubit error at an unknown location.

\subsection*{Status}

Unsolved

\subsection*{Source}

Derived from the gap between the construction and the linear-programming
bound of Rains \sourcecite{ref:5ea7-rains}{Rai99}, recorded as $24$--$26$ in
Table~III of \sourcecite{ref:5ea7-roj}{ROJ19} and in Table~1 of
\sourcecite{ref:5ea7-amh}{AMH26}.  No source located states a conjecture about
the value of $K_{\max}^{(2)}(7,2)$.

\subsection*{Progress}

\begin{itemize}
  \item Rains's Theorem~4 gives pure codes
  \begin{equation}
    ((2m+1,\,3\cdot2^{2m-3},\,2))_2,\qquad m\geq2,
    \label{eq:5ea7-family}
  \end{equation}
  obtained inductively from Lemma~5, which turns a pure $((n,K,2))_2$ code into a
  pure $((n+2,4K,2))_2$ code (numbering of the arXiv version).  Setting $m=3$ in
  Eq.~\eqref{eq:5ea7-family} gives a $((7,24,2))_2$ code; since $24$ is not a
  power of two, this code is non-additive
  \sourcecite{ref:5ea7-rains}{Rai99}.

  \item Rains's Theorem~2 shows that every $((2m+1,K,2))_2$ code satisfies
  \begin{equation}
    K\leq4^{m-1}\Bigl(2-\frac1m\Bigr),
    \label{eq:5ea7-bound}
  \end{equation}
  and that this agrees with the full linear-programming bound.  For seven qubits
  ($m=3$), Eq.~\eqref{eq:5ea7-bound} gives $K\leq80/3$, hence $K\leq26$.  His
  Theorem~3 shows that the bound is not attained at the lengths $2^i+1$ with
  $i\geq3$, where it is an integer, and he remarks that the argument cannot
  strengthen the bound beyond nonattainment.  At seven qubits the bound is not
  an integer, so this argument gives nothing further
  \sourcecite{ref:5ea7-rains}{Rai99}.

  \item The exact rational semidefinite-programming certificates of Angl\`es Munn\'e
  and Huber do not improve this case: their Table~1 (called Table~4.1 in the
  text of Section~4.2) still lists $24$--$26$ for block length $7$ and distance
  $2$ \sourcecite{ref:5ea7-amh}{AMH26}.

  \item Rigby, Olivier, and Jarvis report an exhaustive search over seven-vertex
  graphs up to local complementation and isomorphism: seven classes yield
  distance-two codeword-stabilized (CWS) codes of dimension $24$, and none yields
  dimension $25$ or $26$ (Section~III-A).  A code of dimension $25$ or $26$
  would therefore lie outside the CWS framework
  \sourcecite{ref:5ea7-roj}{ROJ19}.

  \item The two open existence questions are nested.  If $P$ is the projector of a
  $((7,26,2))_2$ code and $P'$ projects onto any $25$-dimensional subspace of its
  range, then $P'=P'P=PP'$ and
  \begin{equation}
    P'EP'=P'(PEP)P'=c_EP'
    \qquad(\operatorname{wt}(E)=1),
    \label{eq:5ea7-subcode}
  \end{equation}
  so $P'$ is a $((7,25,2))_2$ code.  Hence nonexistence at dimension $25$ implies
  $K_{\max}^{(2)}(7,2)=24$, whereas existence at dimension $25$ leaves dimension
  $26$ to decide.  Neither dimension arises as a subcode of a seven-qubit
  distance-two stabilizer code, since such codes have dimension at most $16$
  (Table~II of \sourcecite{ref:5ea7-roj}{ROJ19}).
\end{itemize}

\subsection*{References}

\begin{enumerate}[label={},leftmargin=4.8em,labelsep=0.5em]
  \footnotesize
  \item[\textup{[Rai99]}]\label{ref:5ea7-rains}
  E. M. Rains, ``Quantum codes of minimum distance two,'' \emph{IEEE Transactions on Information Theory} \textbf{45}(1), 266--271 (1999). \href{https://doi.org/10.1109/18.746807}{doi:10.1109/18.746807}; \href{https://arxiv.org/abs/quant-ph/9704043}{arXiv:quant-ph/9704043}.

  \item[\textup{[AMH26]}]\label{ref:5ea7-amh}
  G. Angl\`es Munn\'e and F. Huber, ``SDP bounds on quantum codes: rational certificates,'' arXiv:2603.19901 (2026). \href{https://doi.org/10.48550/arXiv.2603.19901}{doi:10.48550/arXiv.2603.19901}; \href{https://arxiv.org/abs/2603.19901}{arXiv:2603.19901}.

  \item[\textup{[ROJ19]}]\label{ref:5ea7-roj}
  A. Rigby, J. C. Olivier, and P. Jarvis, ``Heuristic construction of codeword stabilized codes,'' \emph{Physical Review A} \textbf{100}, 062303 (2019). \href{https://doi.org/10.1103/PhysRevA.100.062303}{doi:10.1103/PhysRevA.100.062303}; \href{https://arxiv.org/abs/1907.04537}{arXiv:1907.04537}.
\end{enumerate}

\subsection*{Comment}

Literature checked through 15 September 2026: no $((7,25,2))_2$ or
$((7,26,2))_2$ code and no exclusion of either was located, so
$K_{\max}^{(2)}(7,2)\in\{24,25,26\}$ remains open.  The problem is the gap
between an exact construction and the surviving general upper bound; no source
attributes a conjecture $K_{\max}^{(2)}(7,2)=24$.  The error constraints in
Eq.~\eqref{eq:5ea7-kmax} involve only the $21$ weight-one Pauli operators on a
$128$-dimensional space, but they are coupled to the nonlinear condition
$P^2=P$, so counting variables and constraints proves neither existence nor
nonexistence.  By Eq.~\eqref{eq:5ea7-subcode}, a certified exclusion of rank
$25$ would settle the problem, whereas searches confined to CWS codes cannot go
beyond $24$.  The analogous question for nine qubits and distance three is
\href{https://qiqc-op.com/problem/op_0d6ac674da115857/}{the optimal dimension of
nine-qubit distance-three codes}.

\subsection*{Field}

Quantum Error Correction

\subsection*{Topic}

Quantum coding theory

\subsection*{ID}

\texttt{op\_5ea7cb9abb281829}
